\section{INTRODUÇÃO}

Este relatório apresenta a análise e a documentação de uma infraestrutura de rede
corporativa departamentalizada, implementada e simulada na ferramenta Cisco Packet
Tracer (v9.0.1). O cenário exercita a segmentação de domínios de broadcast por meio
de VLANs (\textit{Virtual Local Area Networks}), o transporte de múltiplas VLANs
sobre enlaces de agregação com encapsulamento IEEE 802.1Q, o roteamento inter-VLAN
no modelo \textit{router-on-a-stick} e, na versão final aqui documentada, a
distribuição dinâmica de endereços IP por meio de servidores DHCP hospedados nos
próprios comutadores de acesso.

A relevância do cenário reside na reprodução de um problema recorrente em ambientes
empresariais: setores distintos precisam compartilhar a mesma infraestrutura física de
comutação sem compartilhar o mesmo domínio de broadcast. O isolamento de tráfego entre
departamentos é garantido pelas VLANs, enquanto a comunicação controlada entre eles é
provida pelo roteador central. A atribuição dinâmica de endereços, delegada aos
comutadores, elimina a dependência de um servidor DHCP externo dedicado e distribui
a carga entre os dois equipamentos de acesso.

\subsection{Contextualização e Motivação do Cenário Corporativo}

A arquitetura foi idealizada para atender a uma organização de médio porte dividida
em quatro setores funcionais, cada um mapeado em uma VLAN e em uma sub-rede IP
próprias, acrescida de uma VLAN de gerência/suporte:

\begin{itemize}[leftmargin=2em]
  \item \textbf{VLAN 10 --- VENDAS:} setor comercial. Todos os terminais utilizam
        endereçamento dinâmico via DHCP, refletindo a rotatividade característica
        de equipes de vendas e de dispositivos compartilhados como impressoras.

  \item \textbf{VLAN 20 --- MARKETING:} setor de comunicação e campanhas, com
        endereçamento estático atribuído manualmente às estações.

  \item \textbf{VLAN 30 --- ADM:} setor administrativo. Por tratar dados sensíveis
        de gestão, exige isolamento em relação aos demais domínios de broadcast.

  \item \textbf{VLAN 40 --- TIC:} setor de Tecnologia da Informação e Comunicação,
        responsável pela operação da própria infraestrutura de rede.

  \item \textbf{VLAN 1 --- Gerência/Suporte:} reservada ao tráfego de gerência dos
        ativos de camada 2 por meio de interfaces virtuais de comutação (SVI) e ao
        acesso de laptops de suporte técnico que precisam alcançar e configurar
        qualquer equipamento da rede.
\end{itemize}